\section{What changes and what survives}

The correction changes the interpretation of the threshold results, not the underlying equilibrium construction. This section maps the three propositions back to the original paper so that the scope of the correction is transparent to a reader who is primarily interested in the published results rather than the algebra. Four distinctions are useful. First, the individual trigger associated with Observation 1 cannot be treated as $\gamma$-invariant over asymmetric spillover configurations. The correct object is the zero set $D_i=0$, whose location can vary with $\gamma$. The symmetric formula in Eq. \eqref{eq:symcross} remains useful, but as a special crossing rather than a global threshold characterization. This is the appropriate way to read the reproduced Table 1 benchmarks and the selected comparison in Fig. 1(d): the numerical values need not be discarded, but they cannot carry the general invariance interpretation assigned to them. The correction therefore preserves the calibration while narrowing the inference drawn from it.

Second, the aggregate boundary associated with Observation 4 cannot generally be represented by a scalar threshold in $\beta_1+\beta_2$. At $\lambda=0$, the straight line $\beta_1+\beta_2=1$ intersects the corrected equality set at the symmetric point but fails at regular asymmetric points. Figure \ref{fig:aggregate} visualizes this distinction. More generally, Proposition \ref{prop:aggregate} shows that the corrected boundary can vary with $\gamma$ as well as with the asymmetry of spillovers. Accordingly, the straight boundaries displayed in the published Fig. 3 should be understood as incorrect global representations of the asymmetric equality set. A corrected figure must use $D_A=0$ itself. This is a structural change in the classification geometry even when, for large $\gamma$, the corrected curve lies numerically close to the published line over part of the unit square. Proximity of two curves is not equality of the parameter regions they define.

Third, the fixed average-spillover classification in Observation 5 and the corresponding cutoff equations require qualification. Proposition \ref{prop:misclass} supplies an exact admissible parameter point for which the published cutoff and the actual aggregate R\&D ranking disagree, and continuity makes the disagreement positive-measure. This is a correction of a parameter-region classification. It should not be restated as a reversal of social welfare: the comparison here is between aggregate R\&D levels, the operational object used in the relevant numerical classification. The original paper uses the language of social preference in this part of the numerical analysis, but the mathematical object corrected here is the aggregate R\&D comparison. Nothing in Propositions 1--3 derives a separate welfare function or establishes that a welfare ranking changes at the counterexample. Maintaining this terminology distinction is essential to the narrowness of the note.

Several parts of Ishikawa and Shibata (2021) are unaffected by this correction. The output-stage equilibrium and the closed-form noncooperative and cooperative R\&D solutions are retained. The displayed $\gamma=50$ trigger benchmarks in Table 1 reproduce numerically. The correction also does not overturn the paper's original Propositions 1 and 2 under the dependency considered here. Original Proposition 2 is, in fact, compatible with the more general point that the relevant comparison may depend on $\beta_1,\beta_2,\gamma$, and $\lambda$; what fails later is the reduction of that dependence to fixed scalar cutoffs. Recomputing the branch choice underlying Eq. (47) shifts the classification locally near the corrected aggregate boundary, but the qualitative pattern summarized in Observation 6 survives at the published calibration. Thus Eq. (47) and Fig. 4 are affected through branch selection near the boundary, not because the entire comparative-static pattern must be replaced.

Finally, the paper's broader statement that stronger market competitiveness tends to enlarge the region in which R\&D competition is preferred is not shown here to reverse globally. What changes is the fixed-cutoff rationale used to summarize that relationship over asymmetric spillovers. The corrected geometry is parameter-dependent, so the broad directional interpretation should be stated without the invalid universal threshold. This narrower conclusion is both sufficient and preferable for a correction note: it identifies exactly which mathematical statements require repair while preserving results that do not depend on them. It also explains why the correction can matter editorially even though the equilibrium formulas survive. Readers who use the published cutoffs to classify asymmetric parameter regions can obtain the wrong regime, whereas readers who use the closed-form equilibria directly retain valid inputs. The proper correction is therefore neither to discard the model nor to treat the discrepancy as a typographical detail; it is to replace the threshold summaries by the actual zero sets and to qualify the associated headline interpretation.
For readers, this distinction also determines how the original results should be used. Calculations that substitute parameter values directly into the published equilibrium expressions remain available, subject to the usual regularity conditions. By contrast, applications that infer a regime solely from the reported universal spillover cutoff should recompute the relevant zero-set comparison. The note therefore supplies a repair rule as well as a diagnosis: retain the equilibria, replace the threshold map, and preserve only those qualitative statements that do not require the invalid scalar reduction.

\paragraph{Reproducibility.}
All algebraic identities, exact rational evaluations, and regularity certificates used in Propositions \ref{prop:individual}--\ref{prop:misclass} are verified by deterministic Python/SymPy code. No stochastic simulation is required for any proposition. The repository URL will be inserted at submission: \texttt{[REPOSITORY URL TO BE INSERTED AT SUBMISSION]}.
